\chapter{Introducción}


En los últimos años se han realizado esfuerzos para transformar las carreras que pertenecen a la Escuela de Matemática de la Universidad de Costa Rica. Un paso iniciial se dio con la creación de la carrera de Educación Matemática, que, en el año 2017, abrió las puertas e integró una primera generación de estudiantes. Un año después, se consideró la necesidad de reestructurar las carreras de Matemática y Ciencias Actuariales, para lo cual se inició un proceso de manera conjunta. En el año 2019 se hizo la separación de los procesos con el fin de que se pudieran desarrollar los apartados necesarios para la construcción del perfil de cada carrera por aparte. Este documento constituye una nueva propuesta para el plan de estudios de la carrera de Matemática.

El proceso ha sido largo. En el periodo 2018-2020, se propuso un camino que estuvo mediado por una subcomisión de la Comisión de Docencia de la Escuela de Matemática, en el que se incorporaron las apreciaciones de docentes y estudiantes, así como una revisión de las prácticas llevadas a cabo en las últimas décadas. Esto permitió llegar a establecer el perfil de salida de lo que consideramos será un matemático del siglo XXI \footfullcite{Comision20}. Tal documento se encuentra incluido en el presente trabajo.

Como parte de las incorporaciones de cambio en la reestructuración de la malla curricular, a la luz del trabajo en los marcos referenciales y el establecimiento del perfil de salida, se determinaron cinco áreas curriculares que orientan el trabajo propuesto para alcanzar los propósitos establecidos. Dichas áreas son: habilidades operacionales, pensamiento abstracto, formalismo, pensamiento algorítmico y habilidades investigativas. Este documento no ahonda en tales áreas, pero constituye un fundamento previo para justificar la presencia de estas en la selección y organización de los contenidos, así como la gestión curricular. 

El proceso llevó a la conclusión de la importancia de incorporar dos énfasis en la carrera: Una ruta más teórica o pura que sigue un enfoque más orientado a formar investigadores en áreas puras de la matemática y otro camino más aplicado, dedicado también a formar investigadores pero abierto también a la formación de profesionales que puedan desenvolverse en otros entornos institucionales no académicos o empresariales.

%En este documento se presentan cuatro apartados. Se inicia con la metodología que describe el proceso de los dos años de trabajo para la elaboración del perfil de salida. Se amplía con un contexto teórico y reflexivo de los marcos referenciales que orientan la toma de decisiones para la presentación final de la reestructuración, unida con el aporte de los docentes que participaron en las diferentes actividades donde se recopilaron datos. En los dos últimos apartados se incluyen la declaración de propósitos y el perfil de salida de la carrera de Matemática, aporte que no había sido contemplado anteriormente desde la creación de la carrera, por lo que se considera una transformación en la forma de concebir esta disciplina en el plan de formación. 

Este documento se estructura de acuerdo con los lineamientos establecidos en \cite{Bolanos15} y \cite{CEA24}. 

\section{Metodología}

En esta sección se mencionan algunos aspectos de la metodología que se usó para la conceptualización de la propuesta curricular y la elaboración de este documento.

%se especifican tres aspectos: la conformación del equipo encargado de este documento, los mecanismos que se llevaron a cabo para la obtención de la información, y, por último, una breve noción de lo que representó diseñar este proceso de construcción del perfil de salida, posterior a la elaboración de los marcos referenciales que son pilares fundamentales para esta y otras etapas que se requieren para la reestructuración de la carrera de Matemática. 

\subsection{El equipo de trabajo}

El inicio del perfil de salida viene por mandato de la Vicerrectoría de Docencia (cf.VD-3575-2017) y según designación del Dr. William Ugalde Gómez, director de la Escuela de Matemática, se asignó la responsabilidad. Posterior a una reunión en enero del 2018, junto con las asesoras del CEA, se comprendió la complejidad de esta labor y se conformó un equipo junto con otros tres integrantes, de donde surgió la subcomisión de la Comisión de Docencia que se encargaría de la elaboración del perfil de salida de la carrera. 

El proceso requería una revisión profunda en aspectos epistemológicos y pedagógicos, así como la importancia de dedicar espacios para interactuar con los miembros de la Escuela de Matemática; se acordó en febrero del 2018 la inserción de un estudiante avanzado de la carrera de Enseñanza de la Matemática, quien, posteriormente, se incorporó como miembro de la subcomisión. En agosto del 2019, y durante los siguientes seis meses, se incorporó a otro estudiante con el mismo perfil, que brindó un importante apoyo en la realización de actividades para la obtención de insumos teóricos. Finalmente, la subcomisión concluyó sus labores con tres miembros en julio del 2020 en el contexto del COVID-19, lo cual indujo una participación vía remota y una forma diferente de articular la responsabilidad de concluir el presente documento. 

Los miembros de la subcomisión se encargaron de plantear las estrategias para construir cada uno de los marcos referenciales, orientados por los documentos suministrados y la guía brindada por las asesoras asignadas por el CEA. Además, ellos organizaron las actividades de consulta, realizaron discusión y análisis de la información, y posteriormente, la redacción de cada uno de los apartados. Finalmente, en aquellos aportes que se juzgó necesario, se compartía el trabajo realizado con otros miembros de la Escuela de Matemática, quienes, en muchas ocasiones, dieron importantes contribuciones y una oportuna intervención para asegurar que el trabajo realizado presentaba las características adecuadas. 

A partir del año 20... se cambió la conformación...

Con la llegada del Dr. Pedro Méndez a la dirección de la Escuela de Matemáticas en el año 2025 de nuevo hubo cambios en la conformación ...

\subsection{Las actividades realizadas}

Como parte del proceso se llevaron a cabo numerosas actividades de diversos tipos. Procedemos a describir brevemente las principales.

\subsubsection{Reuniones de la comisión}

...

% \subsection{Espacios de consulta, técnicas de recopilación y análisis de la información}

% Como se observa en la siguiente imagen, se consideraron cinco modalidades para incorporar información pertinente a la redacción del presente documento, cuyo proceso de construcción se justifica según la necesidad de requerir la intervención de otros actores, además de los miembros de la subcomisión. Dicho proceso se prolongó durante dos años y su razón de ser se esboza a continuación. 

\subsubsection{Entrevista}

En un proceso de reflexión inicial para determinar qué es la matemática, se propusieron diferentes áreas de la Matemática, las cuales están bien representadas por miembros de la Escuela de Matemática de la Universidad de Costa Rica; a estas áreas se les llamó identidades. Lo anterior obedece a la complejidad que conlleva establecer este concepto y lo que requiere la construcción de un instrumento denominado V heurística, con el fin de identificar lo teórico y metodológico de cada identidad. Se eligió a uno o dos docentes para que fuesen entrevistados, según la identidad correspondiente. Inicialmente, se determinaron las siguientes identidades: lógica, probabilidad, estadística, geometría, álgebra, análisis, análisis numérico, teoría de números, topología, modelización matemática, análisis funcional y ecuaciones diferenciales.

Con la entrevista se buscaba determinar en qué consiste cada una de las identidades por medio de seis preguntas, cuyos resultados están presentes en el Anexo 1. El proceso de consulta se llevó a cabo entre abril y julio del 2018.

\subsubsection{Grupos focales}

La construcción de cada una de las V heurísticas, según la identidad matemática correspondiente, logró alcanzar el fin esperado excepto en el caso de álgebra, análisis y geometría. Esto puede deberse a que estas tres áreas son las de mayor predominio en la Matemática. Se acordó, por tanto, establecer tres espacios para realizar grupos focales, donde se reunían de cuatro a seis especialistas según la identidad. Los grupos focales procuraban dar respuesta a las preguntas que se diseñaron para la entrevista, con la posibilidad adicional de escuchar otros puntos de vista.

Los tres grupos focales fueron grabados y, después de la transcripción y posterior análisis, se logró construir las tres V heurísticas faltantes. Los grupos focales se aplicaron en el mes de octubre del 2018. 

\subsubsection{Análisis estructurales de la actividad}

Según cada identidad matemática, se procuró un proceso de mayor análisis para profundizar en cada objeto matemático y en el establecimiento de las acciones que lleva a cabo el especialista de esa área. La construcción del análisis estructural de la actividad permitió visualizar la conexión entre el marco socio profesional y el pedagógico, así como la identificación de gran parte de los conocimientos y habilidades descritos en el perfil de salida propuesto. 

La subcomisión diseñó cuatro de estos análisis estructurales; pero, conforme se avanzaba, iban surgiendo más dudas que respuestas de lo que se procuraba alcanzar. Como consecuencia de
muchos momentos de discusión, al considerar el proceso que ya se había desarrollado, tanto en el marco socio profesional como epistemológico, se concluyó que la estadística dejaba de formar parte de las identidades; ya que, esta área está más relacionada con la carrera de ciencias actuariales. Además, durante el desarrollo del análisis estructural del área del análisis matemático, se consideró también que la identidad de análisis funcional se incorporaba a la identidad análisis. En total se llevaron a cabo diez intervenciones para la construcción de los análisis estructurales de la actividad, ya que se acordó la convocatoria de especialistas, según cada área, para diseñar cada uno de estos, con la presencia de al menos un integrante de la subcomisión. Este proceso se extendió alrededor de 8 meses (con la realización de 20 reuniones desde los intentos iniciales, pasando por la incorporación de los especialistas y el análisis de los resultados obtenidos), durante los cuales se alterno con otras actividades del proceso, como lo fue el planeamiento y ejecución de talleres. 


\subsubsection{Talleres}

Como la elaboración de los marcos referenciales es un apartado para construirse según las creencias, vivencias y expectativas de los que conforman la Escuela de Matemática, en cuanto a la carrera de Matemática corresponde, se destinaron tres momentos para la realización de talleres. Estos contaron con la participación de más de veinte personas en cada uno. Para la escogencia de los participantes, se buscó mantener el mejor equilibrio posible en cuanto a género, edad, años de servicio y área de especialización. Además, se solicitó la presencia de algunos estudiantes avanzados de la carrera. 

En el primer taller se extrajo información para iniciar la experiencia de la construcción del marco pedagógico. Se recabó información asociada a las formas en las cuales los participantes enseñan o han aprendido Matemática, las estrategias que aplican actualmente y su visión sobre la futura malla curricular que deberá resultar de este proceso de reestructuración de la carrera. En el segundo taller se trabajó en la obtención de información para el marco epistemológico, para ir construyendo la noción de qué es la matemática, no a partir de lo que dicen las fuentes o especialistas internacionales, sino de la concepción interna de los participantes.

Por último, en el tercer taller, hubo espacio para que los docentes dieran aportes de la posible construcción de una malla curricular, que considerara las áreas curriculares que fueron especificadas previamente. Para esto, se les entregó un material que explicaba las diferentes formas de acceder al aprendizaje según el tipo de curso que se propusiera. En el Anexo 2 se muestra el documento que las asesoras del CEA facilitaron, pero cuyo aporte tendrá mayor incidencia en la propuesta de la reestructuración. Aun así, esto brindó también insumos para la construcción del perfil de salida, e inclusive, para la declaración de propósitos. 

\subsubsection{Consultas, principalmente por correo electrónico}

La subcomisión se reunía con frecuencia para compartir cada uno de los avances e ir estableciendo acuerdos sobre cada aporte generado, en el proceso de construcción del documento. En el transcurso de la experiencia, se vislumbra, por un lado, la oportunidad de enriquecer diferentes apartados al compartir el producto teórico elaborado con otros colegas; y, por el otro, la urgencia de solicitar información ante diferentes escenarios, cuando los propios miembros de la subcomisión consideraban que no tenían la respuesta o alguna forma de dar inicio con esa asignación. La consulta por correo electrónico a diferentes miembros de la Escuela de Matemática fue muy valiosa; se describe, a continuación, en cuáles apartados fue que se contribuyó directamente para la finalización de este documento: principales hitos de la historia de la matemática, historia del desarrollo de la matemática en Costa Rica, prácticas emergentes, decadentes y dominantes, cuadros finales de los análisis estructurales de la actividad, importancia de las identidades matemáticas determinadas en el marco epistemológico, la finalidad de la matemática, perfil de salida, entre otros.
